% results_limitations

\section{Limitations}
\label{sec:limitations}

We state these plainly rather than as hedges: each has a claim ceiling the
results above already respect.

\begin{itemize}

\item \textbf{No users were studied.} Every finding is inferred from public
app-store reviews, a self-selected sample skewed toward strong positive and
strong negative experiences and at best an indirect proxy for lived religious
practice. Claims about user experience here are claims about what reviewers
chose to write down.

\item \textbf{The corpus is an English/US-locale slice, and the applications
were chosen purposively} (\S\ref{sec:method-corpus}), so it under-represents
the regions where most of the world's Muslims live. The aspects most exposed
are the regionally patterned ones --- madhab, calculation method, travel
concession --- so RQ2's finding that madhab complaints are mild and cheap is
not a global claim. Because every application was collected under identical
parameters, this limits generalisation rather than threatening the comparisons
between them.

\item \textbf{Aspect-tagging precision is uneven.} Seven of the 25 aspects with
surviving gold rows fall below 50\%, two at 0\%
(\S\ref{sec:method-measurement}); counts for those seven are paired with
precision-adjusted estimates. The five models draw on aspects at 64--100\%
precision with two exceptions flagged at the point of use. Noise of this kind
attenuates effects toward zero rather than manufacturing them, so it works
against the paper's claims.

\item \textbf{Eleven model terms cannot support inference.} Eight are
non-converged rather than precisely estimated nulls, retained in the
Benjamini--Hochberg denominator of 29, where they are eight of the nine terms
failing correction; the other three return neither a standard error nor a
\(p\)-value and sit outside it (\S\ref{sec:method-models}). Every null reported
here is cited from a fit that converged.

\item \textbf{The topic model characterises a minority of the corpus, and its
structure is seed-sensitive.} 22.7\% of modelled documents fall in the outlier
bucket and 61\% in a single undifferentiated praise cluster; the 39 topics we
label cover about 16\%. The \texttt{rq5\_bloat} sub-model yields 6 to 11 topics
across five fixed seeds (mean pairwise ARI 0.452), so we claim the recurring
themes such a model surfaces, not a specific partition.

\item \textbf{The RQ4 sub-model collapsed around one application's privacy
episode.} Six of ten topics sit on a single theme concerning one application's
user data being sold to US military-linked buyers, holding 58\% of the
sub-model's non-outlier documents. That application is Muslim Pro, whose
November 2020 data-sale reporting we name explicitly because the finding is
specific to that episode rather than a property of the ecosystem; it precludes
further resolution of menstrual-handling discourse (\S\ref{sec:rq4}).

\item \textbf{Illustrative quotations skew toward one application.} Seven of
the 13 reviews quoted come from Athan, the corpus's largest tracker-bearing
application, because quotation followed the volume of available text rather
than a sampling rule --- the same hazard we identify as a methodological
contribution (\S\ref{sec:rq1-coding}). Every claim they illustrate rests on the
corpus-scale or cross-application evidence reported alongside them.

\item \textbf{The affect lexicon is unvalidated.} The guilt and motivation
keyword lists behind RQ1's lexical comparison have no gold sample and no
precision figure, and their vocabulary overlaps ordinary praise and defect
language (\S\ref{sec:method-affect}). RQ1's claim about guilt rests on the
hand-coded passes, which point the same way.

\item \textbf{Two qualitative coding passes, never pooled} (\S\ref{sec:rq1-coding}),
coding 80 reviews between them. They are supporting evidence that sharpens the
corpus-scale results, not a primary contribution.

\item \textbf{Three small-\(N\) case studies are descriptive only.} Inclusion
features ship in 3 of 26 applications; companion hardware and qasr travel
concession each in exactly 1. None supports a between-app statistical estimate,
and we report none.

\item \textbf{Promise-source agreement is weak.} Store descriptions and the
hand-built feature matrix agree at a mean Cohen's \(\kappa\) of 0.272, against
73.1\% mean raw agreement, with four features at \(\kappa = 0\) as a prevalence
artefact (\S\ref{sec:method-promise}). Any claim depending on what an
application promises carries that ceiling.

\item \textbf{The demand-precision sample predates the final extraction.} The
200-row gold sample behind the 91.5\% figure was labelled against an earlier
run; current demand counts come from the regenerated extraction
(\S\ref{sec:method-demand}).

\item \textbf{Topic modelling covers English-identified text only.} Of 324,687
texted reviews, fastText identified 276,422 as English, and only that subset
entered topic modelling, so topic-model findings do not characterise discourse
in the excluded languages.

\item \textbf{Reviews are cross-sectional, not longitudinal.} Nothing links a
user's reviews across time, so we can speak to aggregate patterns but not to
how an individual's relationship with an application changes over their tenure.

\end{itemize}
